% ============================================================================
% 01_introduction.tex
% ============================================================================

\section{Introduction}
\label{sec:introduction}

Buying a first home is among the largest financial decisions most households take, and it carries cultural weight: around 65\% of English households own their home \parencite{ehs2024}, and ownership is widely regarded as financially superior. The public debate on ``buying vs.\ renting'' runs on slogans: renting as ``dead money'' or ``paying your landlord's mortgage'' against weak English house-price returns. Nobody has measured, on monthly data, how a typical first-time buyer fared against an identical renter investing the difference in equities.

This paper runs that race on official English data from January 2005 to June 2026, in a monthly cash-flow-matched simulation after \textcite{beracha2012}. The buyer takes a 10\% deposit and a 30-year term, both at the regulator's median, on a five-year fixed rate at 90\% loan-to-value refixed as the loan amortises, with calibrated maintenance, depreciation, transaction costs and period-accurate stamp duty. The renter occupies the same dwelling at market rent and invests the deposit and purchase costs, plus every monthly difference in outlays, in a global equity index fund (MSCI ACWI, in sterling). I compare the owner's net equity after selling costs with that portfolio.

Three measurement choices shape the results. First, the comparison runs from every monthly entry date, each cohort holding exactly \emph{five years}, roughly a first-time buyer's tenure \parencite{santander_tenure}: 199 cohorts, which isolates \emph{timing}. Second, it is repeated at holding periods from one to ten years. Third, the official England price average reflects the transaction mix while the rent average is dominated by London, so the national price-to-rent ratio (16.9 in June 2026) sits \emph{below every English region} (17.7--22.4) and flatters ownership. I build instead a representative English dwelling, price and rent both population-weighted averages of the nine regions, and run it region by region.

The two strategies were financially comparable, and not only at five years. At five years buying produced greater terminal wealth in 102 of 199 cohorts (51\%), with a median renter/owner wealth ratio of 0.98 and a median gap of \gbp{994} in the buyer's favour. At seven and ten years buying won more often, 62\% and 61\%, but the median renter still finished within 8--11\% of the owner: at ten years the median gap was \gbp{14{,}454} on a dwelling worth \gbp{319{,}488}. Beyond five years buying was ahead more often than not, by margins small enough that renting remained a serious alternative for a household that actually invests the difference every month, which Section~\ref{sec:robustness} shows is where the renter's side of the result comes from.

The distribution behind those medians is asymmetric: buying won more often, renting won bigger. Over five-year holds the mean wealth ratio was 1.22 against a median of 0.98, while the mean gap was \gbp{1{,}744} in the owner's favour. Renting's proportional wins came where the owner's equity had collapsed, so a high ratio sits on a modest sum, while buying won on much larger balance sheets. Renting beat buying twofold in 21 cohorts, threefold in 16 and fourfold in 5, peaking at 4.17 in November 2007: owner \gbp{15{,}769}, renter \gbp{65{,}685}. The ratio's denominator turns negative for some one- and two-year cohorts, so I quote means only from five years out, and never by region, where negative equity leaves a near-zero base.

Timing decides which side of that distribution a household lands on. Renting won every cohort entering in 2006--2009 and 2016; buying won every cohort entering in 2011--2013, 2018 and 2019. The price-to-rent ratio and mortgage rate at purchase jointly account for 56\% of cross-cohort variation, which Section~\ref{sec:predictors} qualifies: the same variables explain 97\% where there is nothing to forecast, so this is largely the user-cost identity; a circular-shift null gives a median $R^2$ of 0.29, leaving $p = 0.034$; and out of sample the model does worse than the historical mean. One regularity survives: renting won all seven cohorts entering above a price-to-rent ratio of 24.4. There is no corresponding floor below which buying was safe: renting also won at the sample's lowest ratio, 17.6. Regional buy-win rates ran from 38\% (North East) to 63\% (West Midlands).

Leverage explains both sides of that record. Global equities returned 10.4\% per year in sterling over the sample against the dwelling's 3.2\%, and the mortgage turns that modest growth into a comparable return on the buyer's own capital; an unlevered cash buyer wins only 8\% of cohorts, and the win rate peaks at 54\% for a 15--20\% deposit against 46\% at 5\%. Leverage also concentrates bad timing: eleven North East and five North West cohorts entering between mid-2007 and mid-2008 ended five years with \emph{negative} net worth after selling costs. The holding-period profile is a feature of this sample rather than of tenure. Buying wins 30\% of one-year cohorts, 51\% at five years and 60--67\% from six to ten; the short end is round-trip transaction costs, roughly thirty per cent of the buyer's initial outlay, which take six years to amortise. Longer holds rest on one crash, one long cheap-credit expansion and one exceptional equity bull market, so I stop at ten years.

Finally, the \textbf{required appreciation rate} (RAR) is the annual house-price growth at which buying exactly breaks even with renting and investing: the user-cost condition of \textcite{poterba1984} and \textcite{himmelberg2005} solved for the capital-gains term. Today's answer turns on what one expects of equities rather than of houses. At June 2026 conditions over five years, the dwelling needs 2.3\% per year to match an equity investor earning a cautious CAPE-implied 5.9\%, but 3.2\% to match one earning the 10.4\% actually delivered. The dwelling's own 3.2\% sits at the top of that range, so the two assumptions bracket the historical outcome, and a single break-even figure would make the case look settled when it is not.

Four contributions follow: the first monthly, cash-flow-matched English backtest with an equity-investing renter, extending the annual, savings-account analysis of \textcite{mostafa2019}, which ends in 2012; the correction to the national price and rent aggregates; the RAR under two openly different views of equity returns; and open code, data and an online implementation.
